\section{Pre-Registration}
\label{supp:prereg}

\textit{Purpose of this supplement.}\;
This supplement summarises the pre-registration scope, amendment
families, and the computational environment used for all training
and evaluation runs.

The factorial study comprised two waves of training
(\S\ref{sec:training}): the
schema-comparison wave fixed schedule and varied schema and encoder;
the configuration wave fixed schema and varied encoder and schedule.
The configuration wave was preceded by a pre-registration that fixed
the factorial design, the primary hypotheses, the schema-selection
rule, the multiple-comparison tier, the bootstrap protocol, and the
matching radius for the rank-inversion analysis. The document
was committed before any configuration-wave run was submitted, so
the analyses reported here followed the plan rather than the data.

Three amendment families are recorded against the lock; the same
three are listed in \S\ref{sec:prereg} of the main text. Amendment~1
(pre-unblinding) refined analysis specifications --- the slope
CI-width reportability gate and the variance-ratio factor set ---
without changing any hypothesis. Amendment~2 handled the
parameter-efficient fine-tuning collapse: the original LoRA arm
was dropped under a pre-committed three-criterion failure rule
rather than re-tuned post-hoc (audit chain in
\S\ref{supp:lora-audit}); the deferred architecture hypothesis
(H5) and the manuscript-preparation reclassification of the
PubMedBERT-base vs BioLinkBERT-base contrast as architecture
rather than scale are recorded under the same amendment family.
Amendment~3 added the inter-annotator audit, the ICC analysis, and
the $\rho$ sensitivity analysis as post-hoc validity checks during
manuscript preparation. None of the three amendments changed a
pre-registered decision rule, threshold, or comparison tier.

The pre-registration document, amendment log, code, and per-run
analysis outputs are openly available; the location is identified in
the Data and Code Availability statement.

\subsection{Computational environment}

Training and evaluation ran on the University of Bristol's BriCS
Isambard-AI system under Slurm, using NVIDIA GH200 Grace Hopper
Superchips (Grace CPU + H100 GPU, 96~GB HBM3 per GPU). Total
compute was approximately 60 GPU-hours for the 310 main training
runs, 8 GPU-hours for evaluation, and 7 GPU-hours for the
falsification audit. The pinned software environment (Python 3.11,
PyTorch 2.6, Transformers 5.4, PEFT 0.19) is recorded in
\code{requirements.lock} at the released code revision.